\chapter*{Ringraziamenti}
\addcontentsline{toc}{chapter}{Ringraziamenti}

A Mamma e Papà, che mi hanno sempre permesso di seguire la mia strada con il massimo sostegno, senza farmi pesare nemmeno una volta tutti i sacrifici che hanno fatto, e continuano a fare, per me.

A Ilaria, la mia sorellona nonché mia chef personale (per sua scelta, chiaramente, ma io di certo non disdegno), da sempre la mia fan numero uno, nonostante odi profondamente la mia incapacità di scrivere in italiano due frasi di senso compiuto di fila.

A J e Day, quelli che considero i miei migliori amici, che c'erano già prima che iniziassi questo tedioso percorso e che mi hanno supportato e, soprattutto, sopportato in ogni singolo momento, anche dopo ogni mio «che fastidio» pronunciato nell'esatto secondo in cui uscivo dall'aula d'esame. (P.S. J mi supporta e sopporta; Day mi supporta e insulta, ma questo è meglio non dirlo ad alta voce\ldots)

A ChainsawSb--- ehm, no, questo non posso dirlo. Meglio nominarvi uno per uno: Luca, Valerio, Frakka e Manolo. A voi, che ho conosciuto in questi tre anni di puro inferno e delirio, con cui ho condiviso alcuni dei momenti più incredibili e surreali concepiti dall'umana stirpe e che, senza esagerare, mi avete permesso di arrivare fino a questo giorno.

A Nonna, Nonno, Asia, Zia e Zio, che chiamo veramente poco\ldots{} No, ok, devo essere realista: che non chiamo praticamente mai, ma che so essere sempre lì, pronti ad aspettare ogni mia notizia, buona o brutta che sia.

E infine a te, Kai. A quando siamo usciti con quel 93 che ci è sembrato una sconfitta. A quando abbiamo fallito quel test d'ingresso e sembravamo destinati a un percorso miserabile.

Eppure eccoci qua, con questa corona in testa, con questo pezzo di carta in mano, consapevoli di non avere nemmeno il tempo per goderceli, perché dobbiamo recuperare le innumerevoli lezioni della magistrale che ci siamo lasciati indietro.

Sì, Kai, magistrale.

Sì, Kai, ce l'abbiamo fatta.
